\chapter{Trading de la volatilité}\label{ch:dv:trading-volatility}

Un trader achète un straddle à la monnaie d'un mois à 18 de volatilité, parce que l'action a réalisé 25 pendant deux mois, et le couvre en delta
chaque jour. Un mois plus tard, l'action a de nouveau réalisé exactement 25, et la position a perdu de l'argent. Pendant trois semaines, l'action a progressé d'un demi-point
de pourcentage par jour, trop doucement pour compenser la décroissance temporelle du straddle alors que son gamma était élevé. Pendant la dernière semaine, elle a bougé de près de 3~\% par jour, mais à ce moment-là elle
était à huit pour cent au-dessus du strike, où le straddle ne conservait presque plus de gamma. Les mêmes 25~\% réalisés, répartis régulièrement sur le mois, auraient
rapporté 3.28 pour une prime de 4.15. Ce chapitre porte sur ce que rapporte réellement une position de volatilité~: le \omterm{def:dv:trading-volatility:scalping}{gamma scalping} et sa dépendance à la trajectoire,
le carry et le roll-down, les trades de skew et de terme, l'attribution quotidienne d'un livre de volatilité, et la prime que les vendeurs de volatilité encaissent.

\section{Gamma scalping}

Le chapitre 4 a établi le \omterm{def:dv:greeks-and-the-hedging-pnl:hpnl}{P\&L de couverture} d'une option couverte en delta~: chaque jour, elle rapporte $\frac12\Gamma S^2\bigl(r^2-\sigma_{\text{imp}}^2\Delta t\bigr)$,
où $r$ est le rendement du jour. Il est positif lorsque le mouvement du jour dépasse la variance implicite, négatif sinon, pondéré par le
\omterm{def:dv:greeks-and-the-hedging-pnl:cashgamma}{cash gamma} de ce jour.

\begin{definition}[Gamma scalping]\label{def:dv:trading-volatility:scalping}
Le \emph{gamma scalping}\index{gamma scalping} consiste à détenir une position longue sur options et à réajuster leur couverture en delta à intervalles réguliers, de sorte que la couverture vende
après les hausses et achète après les baisses~; son P\&L est la somme, sur les intervalles, du \omterm{def:dv:greeks-and-the-hedging-pnl:cashgamma}{cash gamma} multiplié par l'écart entre le carré du rendement réalisé
et la variance implicite de l'intervalle.
\end{definition}

\begin{definition}[Écart implicite-réalisé]\label{def:dv:trading-volatility:spread}
L'\emph{écart implicite-réalisé}\index{ecart@écart implicite-réalisé} d'une position sur option est la différence entre la volatilité implicite à laquelle elle a été négociée et la
volatilité réalisée ensuite par le sous-jacent sur la vie de l'option, grandeur sur laquelle une position couverte en delta parie
et seulement en moyenne.
\end{definition}

La somme est pondérée, et les poids suivent la trajectoire. Le \omterm{def:dv:greeks-and-the-hedging-pnl:cashgamma}{cash gamma} d'un straddle est maximal près du strike et près de l'échéance, et faible loin
du strike. La variance réalisée gagnée là où le gamma est faible compte peu.

\begin{example}[Même volatilité, jours différents]\label{ex:dv:trading-volatility:days}
Un straddle à la monnaie de 21 jours sur une action à 100, acheté à 18 de volatilité (prime 4.15), couvert chaque jour à 18. Trois trajectoires réalisent toutes exactement
25~\%~:
\begin{itemize}
\item \textbf{trajectoire régulière}~: un mouvement de $\pm1.57\%$ chaque jour, à la hausse puis à la baisse. Le straddle rapporte 3.28~;
\item \textbf{les mauvais jours}~: quinze jours de $+0.5\%$ (7.9~\% réalisés), qui portent l'action à 107.8, puis six jours de $\pm2.84\%$ (45~\%
réalisés). Les semaines calmes perdent 0.80 alors que le gamma est élevé~; la semaine agitée, loin du strike, rapporte 0.02. Le total vaut $-0.78$~;
\item \textbf{les bons jours}~: quinze jours de $\pm0.5\%$ au strike, puis la même semaine agitée à cet endroit, avec le gamma le plus élevé du mois. Le
straddle rapporte 2.87.
\end{itemize}
La somme sur le \omterm{def:dv:greeks-and-the-hedging-pnl:cashgamma}{cash gamma} reproduit ces valeurs à quelques dixièmes près (3.05, $-0.86$, 3.03)~: le reste provient des termes d'ordre supérieur des grands mouvements
(la \cref{fig:dv:trading-volatility:scalp}).
\end{example}

\begin{omfigure}
\begin{tikzpicture}
\begin{groupplot}[group style={group size=2 by 1,horizontal sep=1.6cm},
  width=6.6cm,height=5cm,xlabel={jour de cotation},xmin=0,xmax=21,
  tick label style={font=\small},label style={font=\small},grid=major,grid style={gray!25},table/col sep=comma]
\nextgroupplot[ylabel={cours de l'action},ymin=94,ymax=112,
  legend columns=3,legend cell align=left,legend style={font=\footnotesize,at={(1.12,-0.3)},anchor=north,draw=none,fill=none,/tikz/every even column/.append style={column sep=8pt}}]
\addplot[omDef,thick] table[x=day,y=s_uniform]{figdata/derivatives/25-trading-volatility/scalp.csv};
\addplot[omProp,very thick] table[x=day,y=s_wrong]{figdata/derivatives/25-trading-volatility/scalp.csv};
\addplot[omThm,thick,dashed] table[x=day,y=s_right]{figdata/derivatives/25-trading-volatility/scalp.csv};
\addplot[gray,dotted,domain=0:21] {100};
\legend{régulière,mauvais jours,bons jours}
\nextgroupplot[ylabel={P\&L de couverture cumulé},ymin=-2,ymax=4]
\addplot[omDef,thick] table[x=day,y=pnl_uniform]{figdata/derivatives/25-trading-volatility/scalp.csv};
\addplot[omProp,very thick] table[x=day,y=pnl_wrong]{figdata/derivatives/25-trading-volatility/scalp.csv};
\addplot[omThm,thick,dashed] table[x=day,y=pnl_right]{figdata/derivatives/25-trading-volatility/scalp.csv};
\end{groupplot}
\end{tikzpicture}

\omcaption{Un straddle d'un mois acheté à 18 de volatilité et couvert chaque jour le long de trois trajectoires qui réalisent toutes 25~\%. À gauche~: l'action (en pointillé~: le
strike). À droite~: le \omterm{def:dv:greeks-and-the-hedging-pnl:hpnl}{P\&L de couverture} cumulé. Les mauvais jours, le titre quitte le strike sans bruit et bouge violemment là où le straddle n'a presque
plus de gamma restant. Données~: le tutoriel.}\label{fig:dv:trading-volatility:scalp}
\end{omfigure}

Deux conséquences en découlent. Une option couverte en delta ne parie sur l'\omterm{def:dv:trading-volatility:spread}{écart implicite-réalisé} qu'en moyenne sur les trajectoires. Sur une trajectoire donnée, elle parie sur
la variance réalisée pondérée par le gamma, ce que retire un \omterm{def:dv:variance-swaps-and-volatility-derivatives:vs}{swap de variance} (chapitre 14) en maintenant un \omterm{def:dv:greeks-and-the-hedging-pnl:cashgamma}{cash gamma} constant. Et la
fréquence de couverture compte peu pour le P\&L espéré, mais elle compte pour son bruit et pour les coûts de transaction, le compromis que chiffre le chapitre 26.

\section{Carry de volatilité et roll-down}

Une option conservée alors que rien ne bouge gagne ou perd malgré tout~: la décroissance temporelle, et la variation de sa volatilité implicite à mesure que son échéance se raccourcit.

\begin{definition}[Roll-down de volatilité]\label{def:dv:trading-volatility:rolldown}
Le \emph{roll-down de volatilité}\index{roll-down de volatilité} d'une option sur un horizon est la variation de sa volatilité implicite lorsque son temps
jusqu'à l'échéance se raccourcit de l'horizon le long d'une structure par terme inchangée (et, pour une surface en \omterm{def:dv:implied-volatility-and-its-surface:sticky}{sticky delta}, d'un smile inchangé en moneyness)~; avec la
décroissance temporelle, elle constitue le carry de l'option.
\end{definition}

\begin{example}[Roll-down le long d'une courbe ascendante]\label{ex:dv:trading-volatility:carry}
La structure par terme à la monnaie est de 19.70 à un mois, 20.92 à deux, 21.79 à trois et 23.89 à un an. Un straddle de trois mois valant 8.69
détenu pendant un mois, spot et courbe inchangés, vaut 6.81~: un carry de $-1.88$. À volatilité inchangée de 21.79, il aurait perdu 1.59 au titre de la
décroissance temporelle. Les 0.28 restants relèvent du roll-down, la volatilité implicite du straddle baissant de 0.87 point en devenant une option de deux mois. Un calendar (long
le straddle de trois mois, short de 1.73 straddles d'un mois, neutre en \omterm{def:dv:greeks-and-the-hedging-pnl:greeks}{vega}) rapporte à la place 0.13 par jour de carry, payé par un gamma négatif de
$-0.17$ (la \cref{fig:dv:trading-volatility:carry}).
\end{example}

\begin{omfigure}
\begin{tikzpicture}
\begin{groupplot}[group style={group size=2 by 1,horizontal sep=1.6cm},
  width=6.6cm,height=5cm,
  tick label style={font=\small},label style={font=\small},grid=major,grid style={gray!25},table/col sep=comma]
\nextgroupplot[xlabel={échéance (mois)},ylabel={volatilité à la monnaie (\%)},xmin=0,xmax=12,ymin=18,ymax=25]
\addplot[omThm,very thick] table[x=months,y=vol]{figdata/derivatives/25-trading-volatility/term.csv};
\addplot[omProp,only marks,mark=*,mark size=2pt] table[x=months,y=vol]{figdata/derivatives/25-trading-volatility/term_points.csv};
\nextgroupplot[xlabel={jour de cotation},ylabel={valeur du straddle},xmin=0,xmax=21,ymin=6.5,ymax=9,
  legend columns=1,legend cell align=left,legend style={font=\footnotesize,at={(0.5,-0.3)},anchor=north,draw=none,fill=none}]
\addplot[omProp,very thick] table[x=day,y=roll]{figdata/derivatives/25-trading-volatility/carry.csv};
\addplot[omDef,thick,dashed] table[x=day,y=flat]{figdata/derivatives/25-trading-volatility/carry.csv};
\legend{roll-down le long de la courbe,volatilité maintenue à 21.79}
\end{groupplot}
\end{tikzpicture}

\omcaption{À gauche~: une structure par terme à la monnaie ascendante (points à un, deux et trois mois). À droite~: un straddle de trois mois détenu pendant un mois
spot et courbe inchangés~: sa valeur baisse sous l'effet de la décroissance temporelle et du roll-down de sa volatilité. Données~: le
tutoriel.}\label{fig:dv:trading-volatility:carry}
\end{omfigure}

\section{Trades de skew et de structure par terme}

Un calendar spread est un trade de terme~: long de la volatilité d'une échéance, short de celle d'une autre, neutre en \omterm{def:dv:greeks-and-the-hedging-pnl:greeks}{vega} dans l'ensemble. Il profite si la courbe s'aplatit ou s'inverse, et
son carry provient du roll-down. Un risk reversal est un trade de skew~: long d'un call à 25~delta, short d'un put à 25~delta (ou l'inverse), presque
neutre en \omterm{def:dv:greeks-and-the-hedging-pnl:greeks}{vega} et long ou short du skew. Son gain dépend de la manière dont la surface évolue avec le spot, ce qui relève de la règle de marquage du chapitre~7.

\begin{example}[Une mauvaise journée, deux marquages]\label{ex:dv:trading-volatility:rr}
Trois mois avant l'échéance, volatilité à la monnaie de 20~\% et un skew qui place le put à 25~delta (strike 93.6) à 21.3 et le call à 25~delta (106.96) à
18.7. Le desk est short du put, long du call, couvert par 0.5 action. L'action baisse de 5~\% en une journée et le niveau à la monnaie ne bouge pas.
\begin{center}\small
\begin{tabular}{lrrrrr}
\toprule
marquage & gamma & \omterm{def:dv:greeks-and-the-hedging-pnl:greeks}{vega} & \omterm{def:dv:greeks-and-the-hedging-pnl:greeks}{vanna} & non expliqué & total\\
\midrule
\omterm{def:dv:implied-volatility-and-its-surface:sticky}{sticky strike} & 0.053 & $-0.003$ & 0.000 & $-0.077$ & $-0.023$\\
\omterm{def:dv:implied-volatility-and-its-surface:sticky}{sticky delta} & 0.053 & $-0.003$ & 0.112 & $-0.085$ & 0.081\\
\bottomrule
\end{tabular}
\end{center}
Sous \omterm{def:dv:implied-volatility-and-its-surface:sticky}{sticky strike}, chaque option conserve sa volatilité. Sous \omterm{def:dv:implied-volatility-and-its-surface:sticky}{sticky delta}, les volatilités des deux options baissent d'un point, car chaque strike se trouve désormais
plus élevé par rapport au spot et le skew fait baisser la volatilité à mesure que la moneyness augmente~; le \omterm{def:dv:greeks-and-the-hedging-pnl:greeks}{vanna} de la position transforme cela en gain. Le marquage décide du signe de la journée. Les \omterm{def:dv:greeks-and-the-hedging-pnl:greeks}{grecques}
de l'explain laissent 0.08 sans explication, dans un sens comme dans l'autre~: un mouvement de 5~\% est trop grand pour les termes d'ordre deux, et un desk réévalue ces journées-là intégralement.
\end{example}

\section{Attribution du P\&L d'un livre de volatilité}

Chaque jour, le desk explique son P\&L par les \omterm{def:dv:greeks-and-the-hedging-pnl:greeks}{grecques} et compare cette explication à une réévaluation complète. Avec $\delta\sigma_i$ la variation
du marquage propre de l'option $i$,
\[
\text{P\&L}\approx\Delta\,\delta S+\tfrac12\Gamma\,\delta S^2+\Theta\,\delta t+\sum_i\bigl(\mathcal V_i\,\delta\sigma_i+\text{vanna}_i\,\delta S\,\delta\sigma_i
+\tfrac12\,\text{volga}_i\,\delta\sigma_i^2\bigr),
\]
et le reste est le P\&L non expliqué. La construction calcule les deux, par option et par jour. La règle de marquage intervient deux fois~: par
la variation $\delta\sigma_i$ du jour, qui sous \omterm{def:dv:implied-volatility-and-its-surface:sticky}{sticky delta} évolue avec le spot, et par les marquages eux-mêmes.

\begin{example}[Trois mois d'un livre de volatilité]\label{ex:dv:trading-volatility:book}
Un livre long d'un straddle à la monnaie arrivant à échéance dans 0.3 an et short de deux puts 90 de même échéance, couvert en delta chaque jour, sur 63 jours
d'une trajectoire à volatilité stochastique (l'action passe de 100 à 98.6 en repassant par 94.3~; la volatilité de 20 à 19). Totaux, valorisés en \omterm{def:dv:implied-volatility-and-its-surface:sticky}{sticky strike} (\omterm{def:dv:implied-volatility-and-its-surface:sticky}{sticky
delta} entre crochets)~: gamma 2.87 (2.90), \omterm{def:dv:greeks-and-the-hedging-pnl:greeks}{theta} $-1.90$ ($-1.90$), \omterm{def:dv:greeks-and-the-hedging-pnl:greeks}{vega} 1.65 (1.15), \omterm{def:dv:greeks-and-the-hedging-pnl:greeks}{vanna} $-2.24$ ($-1.67$), \omterm{def:dv:greeks-and-the-hedging-pnl:greeks}{volga} $-0.76$ ($-0.62$), non expliqué
$-0.16$ ($-0.09$)~; total $-0.53$ ($-0.22$). Le plus grand montant quotidien non expliqué est de 0.08
(voir la \cref{fig:dv:trading-volatility:explain}).
\end{example}

Les deux marquages s'accordent sur le gamma et le \omterm{def:dv:greeks-and-the-hedging-pnl:greeks}{theta}, qui dépendent du chemin, et divergent sur le \omterm{def:dv:greeks-and-the-hedging-pnl:greeks}{vega} et le \omterm{def:dv:greeks-and-the-hedging-pnl:greeks}{vanna}, qui dépendent de la manière dont la surface est supposée
se déplacer. Elles divergent même sur le total, car le livre est valorisé différemment en cours de route~; l'écart ne disparaît qu'à l'échéance. Un livre
dont le P\&L de \omterm{def:dv:greeks-and-the-hedging-pnl:greeks}{vanna} est aussi important que son P\&L de \omterm{def:dv:greeks-and-the-hedging-pnl:greeks}{vega} est un livre dont la règle de marquage est une position.

\begin{omfigure}
\begin{tikzpicture}
\begin{groupplot}[group style={group size=2 by 1,horizontal sep=1.3cm},
  width=6.6cm,height=5.2cm,xlabel={jour de cotation},xmin=0,xmax=63,ymin=-3.5,ymax=3.5,
  tick label style={font=\small},label style={font=\small},grid=major,grid style={gray!25},table/col sep=comma]
\nextgroupplot[ylabel={P\&L cumulé},title={\small régime sticky strike},
  legend columns=6,legend cell align=left,legend style={font=\footnotesize,at={(1.08,-0.3)},anchor=north,draw=none,fill=none,/tikz/every even column/.append style={column sep=5pt}}]
\addplot[omThm,thick] table[x=day,y=gamma_ss]{figdata/derivatives/25-trading-volatility/explain.csv};
\addplot[omProp,thick] table[x=day,y=theta_ss]{figdata/derivatives/25-trading-volatility/explain.csv};
\addplot[omDef,thick] table[x=day,y=vega_ss]{figdata/derivatives/25-trading-volatility/explain.csv};
\addplot[violet,thick,dashed] table[x=day,y=cross_ss]{figdata/derivatives/25-trading-volatility/explain.csv};
\addplot[gray,thick,dotted] table[x=day,y=unexplained_ss]{figdata/derivatives/25-trading-volatility/explain.csv};
\addplot[black,very thick] table[x=day,y=total_ss]{figdata/derivatives/25-trading-volatility/explain.csv};
\legend{gamma,theta,vega,vanna + volga,non expliqué,total}
\nextgroupplot[title={\small régime sticky delta},yticklabels={}]
\addplot[omThm,thick] table[x=day,y=gamma_sd]{figdata/derivatives/25-trading-volatility/explain.csv};
\addplot[omProp,thick] table[x=day,y=theta_sd]{figdata/derivatives/25-trading-volatility/explain.csv};
\addplot[omDef,thick] table[x=day,y=vega_sd]{figdata/derivatives/25-trading-volatility/explain.csv};
\addplot[violet,thick,dashed] table[x=day,y=cross_sd]{figdata/derivatives/25-trading-volatility/explain.csv};
\addplot[gray,thick,dotted] table[x=day,y=unexplained_sd]{figdata/derivatives/25-trading-volatility/explain.csv};
\addplot[black,very thick] table[x=day,y=total_sd]{figdata/derivatives/25-trading-volatility/explain.csv};
\end{groupplot}
\end{tikzpicture}

\omcaption{Explain quotidien cumulé d'un livre de volatilité couvert en delta (long d'un straddle à la monnaie, short de deux puts 90) sur 63 jours
d'une trajectoire à volatilité stochastique, valorisé en \omterm{def:dv:implied-volatility-and-its-surface:sticky}{sticky strike} (à gauche) et en \omterm{def:dv:implied-volatility-and-its-surface:sticky}{sticky delta} (à droite). Le gamma et le \omterm{def:dv:greeks-and-the-hedging-pnl:greeks}{theta} concordent~; le \omterm{def:dv:greeks-and-the-hedging-pnl:greeks}{vega} et les termes croisés dépendent de la
règle de marquage. Données~: le tutoriel.}\label{fig:dv:trading-volatility:explain}
\end{omfigure}

\section{La prime de risque de variance en pratique}

Le chapitre~14 a défini la \omterm{def:dv:variance-swaps-and-volatility-derivatives:vrp}{prime de risque de variance} comme l'écart entre la variance implicite et la variance réalisée attendue. Bakshi et Kapadia l'ont mesurée à travers
options S\&P~500 couvertes en delta~: les positions couvertes sous-performent zéro, d'autant plus que la volatilité est élevée, ce qui est le signe d'une prime de risque de volatilité de marché négative
pour l'acheteur, et positive pour le vendeur. Israelov et Nielsen décomposent un covered call en ses parts actions, short volatilité
et de retournement des actions. La partie short volatilité présentait un ratio de Sharpe réalisé proche de 1.0, alors qu'elle ne portait
moins d'un dixième du risque de la stratégie.

\begin{dated}{2026-09}{Un indice de référence short volatilité publié}\label{dat:dv:trading-volatility:put}
Le PutWrite Index S\&P~500 de Cboe (PUT) suit une stratégie de short put collatéralisée~: puts S\&P~500 à un mois, à la monnaie, vendus chaque mois, en général le
troisième vendredi, sur un compte investi en bons du Trésor (T-bills) à un et trois mois. Le strike est le strike coté le plus proche, sans lui être supérieur, de la dernière
valeur de l'indice publiée avant 11~h~00\ heure de l'Est (document de méthodologie, consulté le 2026-09).
\end{dated}

La prime est versée en contrepartie du risque de krach, et la distribution des rendements du vendeur le montre~: beaucoup de petits gains, quelques pertes importantes. Le 5 février
2018, le S\&P~500 a baissé de 4~\% tandis que le VIX a bondi de 20 points, le genre de journée qui efface des mois de prime.

\begin{example}[Dix ans de vente de straddles]\label{ex:dv:trading-volatility:vrp}
Chaque mois, un desk vend un straddle à un mois, à la monnaie, à la volatilité que son modèle anticipe, plus deux points, une prime supposée, et
le couvre chaque jour à cette volatilité. L'action suit une trajectoire à volatilité stochastique avec krachs (en moyenne un par an, $-6\%$). Sur 120 mois, la volatilité implicite moyenne vendue
est de 23.5 contre 20.6 réalisée. Le P\&L mensuel moyen représente $+0.44\%$ du spot, avec un écart-type de 1.33 (soit un
ratio de Sharpe de 1.15 par an), 74~\% des mois positifs, un meilleur mois de $+2.89\%$ et un pire de $-5.04\%$~; l'asymétrie est de $-1.35$
(voir la \cref{fig:dv:trading-volatility:vrp}).
\end{example}

\begin{omfigure}
\begin{tikzpicture}
\begin{axis}[width=10.5cm,height=5cm,ybar,bar width=8pt,
  xlabel={P\&L mensuel (\% du spot)},ylabel={mois},
  tick label style={font=\small},label style={font=\small},
  xmin=-6,xmax=3.5,ymin=0,ymax=27,grid=major,grid style={gray!25},
  table/col sep=comma]
\addplot[omProp,fill=omProp!40] table[x=centre,y=count]{figdata/derivatives/25-trading-volatility/shortvol.csv};
\end{axis}
\end{tikzpicture}

\omcaption{P\&L mensuel de la vente d'un straddle à un mois couvert en delta à deux niveaux au-dessus de la volatilité attendue, 120 mois d'un chemin simulé avec
krachs~: la plupart des mois sont gagnants, la queue gauche est longue. Données~: le tutoriel.}\label{fig:dv:trading-volatility:vrp}
\end{omfigure}

\section{Tutoriel~: scalping et explication}\label{tut:dv:trading-volatility:tut}

\begin{tutorial}
\textbf{Objectif.} Faire tourner un straddle couvert en delta sur des chemins construits et comparer le P\&L à la somme des \omterm{def:dv:greeks-and-the-hedging-pnl:cashgamma}{cash gamma}~; détenir un straddle et un calendar sur une
courbe inchangée~; expliquer le mauvais jour d'un risk reversal et les trois mois d'un livre de volatilité sous \omterm{def:dv:implied-volatility-and-its-surface:sticky}{sticky strike} et \omterm{def:dv:implied-volatility-and-its-surface:sticky}{sticky delta}~; faire tourner un programme de vente de
volatilité sur dix ans. \textbf{État final~:} les quatre figures, le tableau du risk reversal et les chiffres du problème du week-end.
\begin{tutsteps}
\item \textbf{L'explication} d'une période, option par option à sa propre valorisation, par rapport à la réévaluation complète~:
\omcode{code/firm/volpnl/firm_volpnl.py}{69}{92}{Explication par les grecques par rapport à la réévaluation complète.}\label{lst:dv:trading-volatility:explain}
\item \textbf{Le scalp}~: couvertures quotidiennes à la volatilité implicite, et la somme des \omterm{def:dv:greeks-and-the-hedging-pnl:cashgamma}{cash gamma}~:
\omcode{code/firm/volpnl/firm_volpnl.py}{106}{119}{Un straddle couvert le long d'un chemin.}\label{lst:dv:trading-volatility:scalp}
\item \textbf{Exécuter} \texttt{dv\_volpnl.scalp()}, \texttt{carry()}, \texttt{risk\_reversal()}, \texttt{book\_explain()}, \texttt{short\_vol()} et
\texttt{fig\_volpnl.py}.
\end{tutsteps}
\textbf{Pour aller plus loin.} Couvrir le straddle des mauvais jours deux fois par jour et à la volatilité réalisée~; ajouter un mouvement de skew à la \omterm{def:dv:implied-volatility-and-its-surface:surface}{surface de volatilité} du livre~; vendre
des puts plutôt que des straddles dans le programme et comparer les queues.
\end{tutorial}

\section{Construction~: P\&L du livre de volatilité}\label{bld:dv:trading-volatility:volpnl}

\begin{build}
\bfield{Rôle} Le P\&L du livre de volatilité de la firme miniature~: positions couvertes réévaluées intégralement, et explication quotidienne par les \omterm{def:dv:greeks-and-the-hedging-pnl:greeks}{grecques} par rapport à celle-ci, selon
l'une ou l'autre des règles de valorisation.
\bfield{Interface} \texttt{Option(strike, expiry, right, qty)}~; \texttt{SkewSurface(atm, skew, ref)} avec \texttt{vol(strike, tau, spot)}
(\texttt{ref=None}~: \omterm{def:dv:implied-volatility-and-its-surface:sticky}{sticky delta}) et \texttt{shifted}~; \texttt{value}, \texttt{book\_greeks}~; \texttt{explain(book, s0, s1, t0, t1, surf0, surf1,
hedge)}~; \texttt{hedged\_pnl(path, book, surfaces)}~; \texttt{gamma\_scalp(returns, strike, days, vol)}~; \texttt{roll\_down(atm, tau, horizon)}.
\bfield{Règles} Chaque option est expliquée à sa propre valorisation~; l'explication est toujours présentée par rapport à la réévaluation complète, jamais à sa place~; la
règle de valorisation est un paramètre explicite.
\bfield{Tests d'acceptation} \texttt{code/firm/volpnl/tests/}~: la part non expliquée est faible pour les petits mouvements et croît comme le cube du mouvement~; la
couverture supprime le P\&L de delta~; les valorisations \omterm{def:dv:implied-volatility-and-its-surface:sticky}{sticky strike} ne bougent pas avec le spot, celles \omterm{def:dv:implied-volatility-and-its-surface:sticky}{sticky delta} si~; un straddle couvert le long d'un chemin réalisant sa
volatilité implicite ne rapporte presque rien~; le roll-down sur des courbes plates et pentues.
\bfield{Pour aller plus loin} \omterm{def:dv:greeks-and-the-hedging-pnl:greeks}{Theta} et roll-down séparés dans l'explication quotidienne~; une explication du \omterm{def:dv:greeks-and-the-hedging-pnl:greeks}{vega} par tranche d'échéance (chapitre 26)~; l'explication des
\omterm{def:dv:autocallables:autocallable}{autocallables} du chapitre 18 à partir de leurs \omterm{def:dv:greeks-and-the-hedging-pnl:greeks}{grecques} obtenues par chocs.
\end{build}

\begin{omsources}
\item G.~Bakshi et N.~Kapadia, «~Delta-hedged gains and the negative market volatility risk premium~», \emph{Review of Financial Studies} 16(2)
(2003) 527--566.
\item R.~Israelov et L.~N.~Nielsen, «~Covered calls uncovered~», \emph{Financial Analysts Journal} 71(6) (2015) 44--57.
\item P.~Carr et L.~Wu, «~Variance risk premiums~», \emph{Review of Financial Studies} 22(3) (2009) 1311--1341.
\item Cboe Global Indices, «~Cboe S\&P 500 PutWrite Indices methodology~» (consulté le 2026-09).
\item Bank for International Settlements, \emph{Quarterly Review}, mars 2018, encadré sur la turbulence des marchés actions du 5 février.
\end{omsources}

\section{Exercices}

\begin{exercise}[$\star$]\label{exo:dv:trading-volatility:1}
Un straddle à un mois à la monnaie sur une action à 100 a un \omterm{def:dv:greeks-and-the-hedging-pnl:cashgamma}{cash gamma} d'environ 770 à l'origine. Que rapporte une journée avec un mouvement de 2~\% à une
volatilité implicite de 18~? Et une journée sans mouvement~?
\end{exercise}

\begin{exercise}[$\star$]\label{exo:dv:trading-volatility:2}
Pourquoi un straddle acheté perd-il de l'argent sur le chemin des mauvais jours alors que la volatilité réalisée dépasse l'implicite de sept points~?
\end{exercise}

\begin{exercise}[$\star$]\label{exo:dv:trading-volatility:3}
Décomposer le carry sur un mois d'un straddle à trois mois entre \omterm{def:dv:greeks-and-the-hedging-pnl:greeks}{theta} et roll-down, et dire ce qu'une courbe inversée changerait.
\end{exercise}

\begin{exercise}[$\star\star$]\label{exo:dv:trading-volatility:4}
Expliquer pourquoi le calendar dégage un carry positif et ce qu'il paie en contrepartie.
\end{exercise}

\begin{exercise}[$\star\star$]\label{exo:dv:trading-volatility:5}
Montrer que sous \omterm{def:dv:implied-volatility-and-its-surface:sticky}{sticky delta}, le niveau à la monnaie étant fixé, la volatilité de chaque strike varie de $-s\,\ln(S_1/S_0)/\sqrt\tau$ lorsque le spot
passe de $S_0$ à $S_1$ (skew $s$), et calculer cette variation pour la journée du risk reversal.
\end{exercise}

\begin{exercise}[$\star\star$]\label{exo:dv:trading-volatility:6}
Pourquoi gamma et \omterm{def:dv:greeks-and-the-hedging-pnl:greeks}{theta} coïncident-ils sous les deux règles de valorisation dans l'explication du livre, alors que \omterm{def:dv:greeks-and-the-hedging-pnl:greeks}{vega} et \omterm{def:dv:greeks-and-the-hedging-pnl:greeks}{vanna} ne coïncident pas~?
\end{exercise}

\begin{exercise}[$\star\star\star$]\label{exo:dv:trading-volatility:7}
\emph{Programmation.} Relancer le programme de vente de volatilité sans prime. Quels sont le P\&L moyen et le ratio de Sharpe~? Le P\&L moyen est-il
significativement différent de zéro, et comment la volatilité implicite vendue peut-elle dépasser la volatilité réalisée moyenne sans prime~?
\end{exercise}

\begin{exercise}[$\star\star\star$]\label{exo:dv:trading-volatility:8}
\emph{Trouver l'erreur.} «~Notre livre de vente de volatilité affiche un ratio de Sharpe supérieur à 1 et trois mois gagnants sur quatre~; son risque est faible.~»
\end{exercise}

\section{Problème~: Les mauvais jours}

\begin{problem}[{Problème du week-end --- quand 25 réalisé perd face à 18 implicite}]\label{pb:dv:trading-volatility:1}
Le trader de l'exemple introductif a acheté un straddle à 21 jours à la monnaie sur une action à 100 avec une volatilité de 18 et l'a couvert chaque jour. La volatilité réalisée de l'action est de 25~\%.

\medskip\noindent\textbf{Partie I --- La transaction.}
\begin{enumerate}
\item Combien coûtait le straddle, et sur quoi le trader pariait-il~?
\item Donner le mouvement quotidien qui réalise 25~\% de manière régulière, et le P\&L du chemin régulier.
\item Donner le \omterm{def:dv:greeks-and-the-hedging-pnl:cashgamma}{cash gamma} quotidien à l'origine, et le mouvement quotidien d'équilibre à la volatilité implicite de 18.
\item De combien le P\&L du chemin régulier diffère-t-il de la somme des \omterm{def:dv:greeks-and-the-hedging-pnl:cashgamma}{cash gamma}, et pourquoi~?
\item Combien un \omterm{def:dv:variance-swaps-and-volatility-derivatives:vs}{swap de variance} sur le même notionnel aurait-il rapporté sur l'un quelconque des trois chemins~?
\end{enumerate}

\medskip\noindent\textbf{Partie II --- Les mauvais jours.}
\begin{enumerate}[resume]
\item Décrire le chemin des mauvais jours et vérifier qu'il réalise 25~\%.
\item Décomposer son P\&L entre les semaines calmes et la semaine agitée.
\item Où se trouvait le spot pendant la semaine agitée, et quel était le gamma du straddle à ce moment par rapport à l'origine~?
\item Que rapporte le chemin des bons jours, et pourquoi moins que le chemin régulier~?
\item Que dit la comparaison de l'\omterm{def:dv:trading-volatility:spread}{écart implicite-réalisé} comme signal de trading~?
\end{enumerate}

\medskip\noindent\textbf{Partie III --- Autres façons d'être acheteur de volatilité.}
\begin{enumerate}[resume]
\item Comment le recalage du strike, lorsque le spot a dérivé, aurait-il modifié le résultat~?
\item Quel est le carry de la détention d'un straddle à trois mois pendant un mois sur la courbe du chapitre, et quelle part en est due au roll-down~?
\item Que rapporte par jour un calendar neutre en \omterm{def:dv:greeks-and-the-hedging-pnl:greeks}{vega}, et quel risque prend-il~?
\item Sur une baisse de 5~\%, que rapporte un risk reversal (vente de put, achat de call) sous \omterm{def:dv:implied-volatility-and-its-surface:sticky}{sticky strike} et sous \omterm{def:dv:implied-volatility-and-its-surface:sticky}{sticky delta}~?
\item Pourquoi un desk doit-il réévaluer une telle journée intégralement plutôt que par les \omterm{def:dv:greeks-and-the-hedging-pnl:greeks}{grecques}~?
\end{enumerate}

\medskip\noindent\textbf{Partie IV --- Jugement.}
\begin{enumerate}[resume]
\item La trader s'est-il trompé sur la volatilité~?
\item Que rapporteriez-vous dans l'attribution du P\&L pour le mois~?
\item Comment la distribution du programme de vente de volatilité explique-t-elle que les vendeurs soient rémunérés~?
\item Énoncer le \emph{résultat nommé}~: le P\&L gamma d'un straddle acheteur lorsque la même volatilité réalisée se produit avec des mouvements concentrés sur des jours à faible gamma,
contre des mouvements répartis uniformément.
\item En une phrase~: que rapporte un produit optionnel couvert en delta~?
\end{enumerate}
\end{problem}

\section{Questions d'entretien}

\begin{interviewq}[$\star$ \iqroles{trader}]\label{iq:dv:trading-volatility:1}
Vous êtes acheteur d'un straddle couvert en delta. D'où provient votre P\&L~?
\end{interviewq}

\begin{interviewq}[$\star\star$ \iqroles{trader, chercheur}]\label{iq:dv:trading-volatility:2}
La volatilité réalisée a dépassé la volatilité implicite et votre straddle couvert a perdu de l'argent. Comment cela s'explique-t-il~?
\end{interviewq}

\begin{interviewq}[$\star\star$ \iqroles{trader}]\label{iq:dv:trading-volatility:3}
Quel est le carry d'un calendar spread sur une structure par terme à pente ascendante~?
\end{interviewq}

\begin{interviewq}[$\star\star$ \iqroles{risque, développeur}]\label{iq:dv:trading-volatility:4}
Comment construiriez-vous un explain quotidien du P\&L pour un livre d'options, et que faites-vous de la partie non expliquée~?
\end{interviewq}

\begin{interviewq}[$\star\star$ \iqroles{trader, risque}]\label{iq:dv:trading-volatility:5}
Le même livre affiche un P\&L de \omterm{def:dv:greeks-and-the-hedging-pnl:greeks}{vega} différent sous \omterm{def:dv:implied-volatility-and-its-surface:sticky}{sticky strike} et sous \omterm{def:dv:implied-volatility-and-its-surface:sticky}{sticky delta}. Lequel est juste~?
\end{interviewq}

\begin{interviewq}[$\star\star\star$ \iqroles{chercheur}]\label{iq:dv:trading-volatility:6}
La \omterm{def:dv:variance-swaps-and-volatility-derivatives:vrp}{prime de risque de variance} est-elle un repas gratuit~? Comment dimensionneriez-vous une stratégie de vente de volatilité~?
\end{interviewq}
